\subsection{An antisymmetric residue criterion for tangency}\label{sec:zero-tangency}

The initial horizontal motion of a simple critical-line zero can be separated from the symmetric part of a deformation.

\begin{proposition}\label{prop:zero-tangency-residue}
Suppose $\Lambda_0(1-s)=\Lambda_0(s)$ and $\Lambda_0(\bar s)=\overline{\Lambda_0(s)}$. Let $\rho$ be a simple zero with $\Re\rho=1/2$, and let
\[
 \Lambda_t(s)=\Lambda_0(s)+t\Phi(s)+O(t^2)
\]
be a holomorphic deformation near $\rho$ and $1-\rho$, with $\Phi$ regular there and satisfying Schwarz symmetry. Define
\[
 \Phi_-(s)=\frac{\Phi(s)-\Phi(1-s)}2.
\]
For the simple-zero branch through $\rho$,
\[
 \Re\rho'(0)=-\operatorname*{Res}_{s=\rho}
                   \frac{\Phi_-(s)}{\Lambda_0(s)}.
\]
Consequently the initial velocities at a collection of simple critical-line zeros are all tangent to that line if and only if $\Phi_-/\Lambda_0$ is regular at every zero in that collection.
\end{proposition}

\begin{proof}
The implicit-function theorem gives $\rho'(0)=-\Phi(\rho)/\Lambda_0'(\rho)$. Since $1-\rho=\bar\rho$, differentiated reflection and Schwarz symmetry imply
\[
 \overline{\Lambda_0'(\rho)}=-\Lambda_0'(\rho).
\]
This derivative is nonzero and purely imaginary. Also $\Phi_-(\rho)=i\Im\Phi(\rho)$. Hence the real part of $\Phi(\rho)/\Lambda_0'(\rho)$ equals $\Phi_-(\rho)/\Lambda_0'(\rho)$, which is the stated residue. A holomorphic numerator divided by a function with a simple zero has a removable singularity exactly when that residue vanishes. This proves the equivalence.
\end{proof}

For the Hurwitz family, put $C(s)=\pi^{-s/2}\Gamma(s/2)$ and $\Lambda_t(s)=C(s)\zeta(s,1+t)$. The completed reflection formula and the Hurwitz parameter derivative are the classical identities in \citet[equations~25.4.3--25.4.4]{DLMFReflection} and \citet[equation~25.11.17]{DLMFShift}. At a nontrivial simple critical-line zero,
\[
 \Phi(s)=-sC(s)\zeta(s+1),\qquad
 \chi(s)=\frac{C(1-s)}{C(s)}
 =\pi^{s-1/2}\frac{\Gamma((1-s)/2)}{\Gamma(s/2)}.
\]
Thus, defining
\begin{equation}\label{eq:hurwitz-antisymmetric-G}
 G(s)=\frac{s\zeta(s+1)-(1-s)\chi(s)\zeta(2-s)}{2\zeta(s)},
\end{equation}
one has
\[
 G=-\frac{\Phi_-}{\Lambda_0},\qquad
 G(1-s)=-G(s),\qquad
 \Re\rho'(0)=\operatorname*{Res}_{s=\rho}G(s).
\]
The sign in the last formula is positive because $G$ is the negative of the quotient in Proposition~\ref{prop:zero-tangency-residue}. These are meromorphic identities and a first-order tangency test. They do not assert that a tangent branch remains on the critical line at higher orders.
